% sec-02: the queued orders and the single approval step.  Table 20 and the
% actionbox.
\section{The Queue}

Six orders are sized, priced by basis rather than by figure, and present in no
broker system. Every one is entered on the next open market session after the
release approval, and nowhere today.

\begin{apptable}
\begin{tabularx}{\textwidth}{>{\raggedright\arraybackslash}p{0.5cm} >{\raggedright\arraybackslash}p{3.6cm} >{\raggedright\arraybackslash}p{1.3cm} >{\raggedright\arraybackslash}p{2.6cm} >{\raggedright\arraybackslash}X}
\headrow \thead{\#} & \thead{Instrument} & \thead{Side} & \thead{Type} & \thead{Limit basis, size, time in force}\\
\midrule
1 & U.S. Treasury bill, about 3 months & Buy & Non-competitive at auction & Not applicable; 20 pct of reserve; auction \\
2 & U.S. Treasury bill, about 6 months & Buy & Non-competitive at auction & Not applicable; 20 pct; auction \\
3 & U.S. Treasury bill, about 9 months & Buy & Non-competitive at auction & Not applicable; 20 pct; auction \\
4 & U.S. Treasury note or bill, about 12 months & Buy & Limit, secondary market & Last session's close plus one tick; 20 pct; day \\
5 & Short-duration Treasury ETF, one fund only & Buy & Limit, day, never market & Last session's close plus one cent; 15 pct; day \\
6 & Government money market sweep & Hold & Sweep & Not applicable; 5 pct residual; continuous \\
\end{tabularx}
\end{apptable}
\vspace{-0.60cm}
\tabcap{Table 20. Six queued orders with instrument, side, type, and the basis\\
each limit is set from, since a price written today would be three days stale.}

\subsection{Why a basis and not a price}

A limit written before a closed weekend and a holiday is three days stale by the
time the market opens. So the queue records the basis, which is the last
completed session's close plus one increment, and the actual figure is written
into the broker instruction on the morning of entry. The same rule governed the
day 1 instruction and is restated here because the interval is longer.

Time in force is **day** on every order that has one, and never good-till-
canceled. A standing order written a session in advance can fill days later at a
price the instruction never contemplated, which converts a queued order into a
standing exposure \cite{treasuryauctions2026}.

\subsection{The auction re-check, which is not optional}

Treasury auctions do not settle on a federal holiday, and an auction calendar
read before one will be wrong for the following week. Before any of orders 1 to
3 is placed, the auction announcements page is re-read and three things are
confirmed: the next auction date for each tenor, its settlement date, and the
broker's own cutoff relative to that auction. A non-competitive bid submitted
after a broker's internal cutoff is not submitted at all.

\subsection{The branch these orders sit under}

No financing instrument has been selected, so every order above sits under the
unfinanced nine-month branch established on day 1. If an instrument is selected
before the next session, the queue is re-cut against the twenty-four-month branch
before it is entered, rather than entered and then amended.

\subsection{What is not queued, and why the omissions are deliberate}

The site start-up pool of the preceding day is not in Table 20. It is held in a
government sweep, one short bill, and one liquid fund, and its draws release
against artifacts produced by an institution rather than against a market
session. Queuing it here would imply a schedule it does not have and would
suggest to a reader that money is committed when none is.

Subscription funds are not queued because none exist and no offering exists.
Federal award funds are not queued because none exist, and would not be queued in
any case: they are drawn under the award's own terms and are never commingled
with the operating reserve in either direction.

Three sources of money, three treatments, and the separation is established
before the first dollar rather than reconstructed afterward.

\subsection{The six checks that stand between this table and an entered order}

The queue is not self-executing and the approval below does not make it so. Six
checks run on the next session before any order is transmitted, and each of them
has failed for somebody at some point, which is why each is written down rather
than assumed.

Is the market actually open, confirmed against the exchange calendar rather than
against an assumption about which Monday it is. Has the release approval been
given. Has the financing decision changed the branch, in which case the queue is
re-cut before entry rather than entered and amended. Has the auction calendar
been re-read, which is the check most likely to be skipped and most likely to
matter. Is the limit price written this morning rather than carried from this
file. And is the reserve balance current, since every size above is a share of it.

\subsection{The single approval step}

\actionbox{One decision is open on this day}{Approve the release list in Table 16
and the queue in Table 20, so that the next open session begins with both already
authorized. One approval on a closed day replaces five separate approvals taken
under time pressure on the following session. Nothing is sent or entered by this
approval: the release itself happens on the next session, against the six checks
in \appfile{funding/auto-fund/07Sep26/investing}.}
